\documentclass[11pt]{article}

% ---------- Packages ----------
\usepackage[a4paper,margin=2.5cm]{geometry}
\usepackage[utf8]{inputenc}
\usepackage[T1]{fontenc}
\usepackage{amsmath}
\usepackage{graphicx}
\usepackage{xcolor}
\usepackage{tabularx}
\usepackage{booktabs}
\usepackage{enumitem}
\usepackage{hyperref}
\usepackage[explicit]{titlesec}
\usepackage{fancyhdr}
\usepackage{caption}

% ---------- Colors ----------
\definecolor{etsitorange}{RGB}{237,125,49}
\definecolor{etsitblue}{RGB}{31,73,125}

% ---------- Section banners (orange bar, like the template) ----------
\titleformat{\section}
  {\bfseries\Large\color{white}}
  {}{0pt}
  {\colorbox{etsitorange}{\parbox{\dimexpr\textwidth-4\fboxsep-12pt\relax}{\centering\thesection.\ #1}}}
\titlespacing*{\section}{0pt}{18pt}{12pt}

% Unnumbered \section* (used internally by \tableofcontents, \listoftables,
% \listoffigures) gets the same orange banner, without a number.
\titleformat{name=\section,numberless}
  {\bfseries\Large\color{white}}
  {}{0pt}
  {\colorbox{etsitorange}{\parbox{\dimexpr\textwidth-4\fboxsep-12pt\relax}{\centering#1}}}
\titlespacing*{name=\section,numberless}{0pt}{18pt}{12pt}

\titleformat{\subsection}
  {\bfseries\large}
  {\thesubsection.}{0.5em}{#1}
\titlespacing*{\subsection}{0pt}{12pt}{6pt}

% Appendix banner (no numbering, e.g. "APPENDIX A: ...")
\newcommand{\appendixbanner}[1]{%
  \colorbox{etsitorange}{\parbox{\dimexpr\textwidth-4\fboxsep-12pt\relax}{\centering\bfseries\Large\color{white} #1}}%
  \vspace{12pt}
}

\renewcommand{\contentsname}{CONTENTS}
\renewcommand{\listtablename}{LIST OF TABLES}
\renewcommand{\listfigurename}{LIST OF FIGURES}

\pagestyle{fancy}
\fancyhf{}
\renewcommand{\headrulewidth}{0pt}
\fancyfoot[C]{\thepage}

\usepackage{float}
\usepackage{array}
\usepackage{cite}
\hypersetup{hidelinks}
\graphicspath{{figures/}}
\setlength{\emergencystretch}{2em}  % avoid long \texttt paths running into the margin

% Numbers quoted in the text, generated from results/logs/*.json by report/make_numbers.py
\input{numbers}

% Claims still to be checked against a quoted source (redefine as empty for the final version)
\newcommand{\todoverify}[1]{\textcolor{red}{\,\textsf{\footnotesize[verify: #1]}}}

\begin{document}

% =====================================================================
% COVER PAGE
% =====================================================================
\begin{titlepage}
\centering
\vspace*{1cm}
\includegraphics[width=0.4\textwidth]{figures/etsit_logo.jpeg}\par
\vspace{1.5cm}

{\color{etsitblue}\bfseries\Large Course:}\par
\vspace{0.5cm}
{\bfseries\huge Predictive and Descriptive Learning}\par
\vspace{1cm}
{\Large Half-term Activities Report}\par
{\bfseries\Large Developing Machine Learning models}\par
\vspace{0.5cm}
\rule{0.7\textwidth}{0.4pt}\par
\vspace{3cm}

\large Name: \underline{\hspace{2mm}Oriana Andrea Ferraz De Jesus\hspace{2mm}}\par
\vspace{0.5cm}
\hfill Date: \underline{\hspace{2mm}October 2026\hspace{2mm}}\par

\end{titlepage}

% =====================================================================
% TABLE OF CONTENTS / LIST OF TABLES / LIST OF FIGURES
% =====================================================================
\tableofcontents
\vspace{1cm}

\listoftables
\vspace{1cm}

\listoffigures

\newpage

% =====================================================================
% 1. DEVELOPING MACHINE LEARNING MODELS
% =====================================================================
\section{Developing Machine Learning models}

\subsection{Problem description}

Diffuse gliomas are primary tumours of the central nervous system, and their standard-of-care
treatment ``is predicated on tumor grade'' \cite{tasci2022}. For decades they were graded from
histology alone, but the histological classification of
lower-grade gliomas ``suffers from high intraobserver and interobserver variability and does not
adequately predict clinical outcomes'' \cite{tcga2015lgg}. The 2016 and 2021 World Health
Organization (WHO) classifications therefore integrate molecular markers with histology
\cite{louis2016,louis2021}: in the 2021 scheme, the common adult-type diffuse gliomas comprise ``only 3 types'':
astrocytoma, IDH-mutant (CNS WHO grades 2--4), oligodendroglioma, IDH-mutant and
1p/19q-codeleted (grades 2--3), and glioblastoma, IDH-wildtype (grade 4) \cite{louis2021}.
The grade matters because prognosis differs widely: in the TCGA lower-grade glioma cohort,
patients with an IDH mutation and 1p/19q codeletion had a median survival of 8.0 years, and those
with an IDH mutation without codeletion of 6.3 years \cite{tcga2015lgg}, whereas glioblastoma,
which accounts for 14.3\% of all primary brain and other CNS tumours in the United States and is
diagnosed at a median age of 65 years, has a five-year relative survival of only 6.8\%
\cite{ostrom2021}.

The molecular markers relevant to this project are point mutations. Gliomas without IDH1 or
IDH2 mutations tend to carry other alterations: among anaplastic astrocytomas and glioblastomas
with wild-type IDH genes, \cite{yan2009} found few TP53 mutations (18\%) and frequent alterations
of PTEN, EGFR, CDKN2A or CDKN2B (74\%). In routine practice these alterations are detected with
immunohistochemistry, sequencing and copy-number assays; for example, fluorescence \emph{in situ}
hybridisation detects gains of chromosome 7 and losses of chromosome 10, the 1p/19q codeletion,
homozygous deletions such as CDKN2A/B, and gene amplifications \cite{sahm2023}.

This motivates the questions of this project, stated as hypotheses before any model was trained:
\begin{enumerate}[label=\textbf{H\arabic*.}, leftmargin=*]
  \item The mutation status of 20 genes, together with age, sex and race, contains enough
        information to separate lower-grade glioma (LGG) from glioblastoma (GBM) well above a trivial
        classifier, and simple interpretable models perform close to complex ones.
  \item Without using the grade, the same profile explains part of the variability of the age at
        diagnosis, mostly through the markers that also define the grade.
  \item Without any label, the mutation profiles form groups that are related to, but not
        identical with, the grade.
\end{enumerate}
After the analysis of H1--H3, two complementary questions were added, and they are reported as such
rather than as hypotheses fixed in advance:
\begin{enumerate}[label=\textbf{Q\arabic*.}, leftmargin=*, start=4]
  \item Can the same profile go beyond the grade and recognise the histological family
        (astrocytoma, oligodendroglioma, mixed glioma, glioblastoma)?
  \item Under explicit causal assumptions, how much of the association between an IDH mutation and a
        younger age at diagnosis goes through the grade?
\end{enumerate}
The objective is not to beat published accuracies but to apply and critically compare the
models of the course under a validation protocol that gives honest estimates.

\subsection{Machine Learning approach}

The three hypotheses map onto the three learning tasks required in the course:
\begin{itemize}
  \item \textbf{Classification} (H1): target \texttt{Grade} (LGG = 0, GBM = 1), predictors age,
        sex, race and 20 mutations. Models: penalised logistic regression, naive Bayes, a pruned
        classification tree, random forest, gradient boosting (XGBoost) and support vector machines.
  \item \textbf{Regression} (H2): target \texttt{Age\_\allowbreak{}at\_\allowbreak{}diagnosis} (years), predictors sex, race
        and 20 mutations, deliberately \emph{without} the grade. Models: multiple linear regression,
        Ridge, Lasso, a regression tree, random forest, XGBoost and support vector regression.
  \item \textbf{Multiclass classification} (Q4): target the histological family derived from the
        diagnosis (four classes), same predictors as H1 and \emph{without} the grade. Models:
        multinomial logistic regression, a pruned tree, random forest and XGBoost. The winning models of
        both classification tasks are then explained with SHAP values.
  \item \textbf{Causal analysis} (Q5): a causal diagram and a mediation decomposition with linear
        regression, with the grade as mediator between IDH and the age at diagnosis.
  \item \textbf{Unsupervised learning} (H3): PCA, multidimensional scaling and t-SNE with a
        Jaccard distance, and $k$-means and hierarchical clustering of the 20 mutation indicators;
        the grade is used only afterwards, to describe the clusters.
\end{itemize}
Every supervised model is evaluated with the same nested cross-validation and compared with a
trivial baseline (Section~\ref{sec:prdl-setup}); the tools and the implementation are described
in the companion MLLB report.

% =====================================================================
% 2. DATA DESCRIPTION
% =====================================================================
\section{Data description}

\subsection{Data bases}

The data are the \emph{Glioma Grading Clinical and Mutation Features} dataset of the UCI Machine
Learning Repository \cite{tasci2022dataset}, built from the TCGA lower-grade glioma and
glioblastoma projects \cite{tcga2015lgg,brennan2013} and introduced in \cite{tasci2022}. The original
export contains \NRaw{} patients; the encoded UCI version contains \NClean{}. Each patient is
described by age at diagnosis, sex, race, the somatic mutation status (mutated or not) of 20 genes,
the TCGA project (TCGA-LGG or TCGA-GBM), the histological diagnosis and the grade. The final dataset
has \NLGG{} LGG and \NGBM{} GBM patients (\PctGBM\% GBM).

Three properties of these data limit what any model can learn, and they shape the critical analysis
of Section~\ref{sec:prdl-results}:
\begin{itemize}
  \item \textbf{Labels from older criteria.} The cohorts were collected over many years
        (TCGA-GBM between 1989 and 2011, with 76\% of patients diagnosed in or after 2002
        \cite{zakharova2023}) and graded before the 2021 WHO classification. TCGA glioma profiles
        have since been re-annotated following the 2021 criteria \cite{zakharova2023,leiriademendonca2025}%
        \todoverify{quote from 41.pdf or 42.pdf stating how many TCGA cases change category when the
        original labels are replaced by WHO 2021 labels}.
        The grade used here as ground truth is therefore partly outdated: under the current criteria,
        some tumours labelled LGG would be classified differently.
  \item \textbf{Point mutations only.} Several markers that now define the type or the grade are
        copy-number or promoter alterations absent from the dataset: the 1p/19q codeletion
        (oligodendroglioma) \cite{louis2021}; the homozygous deletion of CDKN2A/B, ``a molecular marker
        of WHO grade 4 in IDH-mutant astrocytomas''; and the TERT promoter mutation, EGFR amplification
        or the +7/$-$10 signature, which define glioblastoma \cite{weller2021}. These alterations are
        detected with copy-number assays such as FISH \cite{sahm2023}. A model can only approximate them
        through correlated point mutations: in the TCGA lower-grade gliomas, the 1p/19q-codeleted tumours
        ``harbored mutations in CIC, FUBP1, NOTCH1, and the TERT promoter'' \cite{tcga2015lgg}.
  \item \textbf{A homogeneous cohort.} \PctWhite\% of the patients are white (\NWhite{} of
        \NClean); only \NBlack{} are Black or African American, \NAsian{} Asian and \NNative{}
        American Indian or Alaska Native. Any conclusion about race, and the transfer of the models
        to other populations, is unsupported.
\end{itemize}

\subsection{Data preparation}

The preparation followed three principles, each of which can be justified from theory.

\textbf{Reproducibility of the published cleaning.} Rather than loading the UCI file, the clean
dataset was rebuilt from the original export, keeping the patient identifier and the diagnosis,
and checked value by value against the UCI file. The \NRemoved{} excluded patients
(\NRemovedAllMissing{} with every clinical field missing, \NRemovedRace{} with race not reported and
\NRemovedAge{} with missing age) are split almost evenly between classes (\NRemovedLGG{} LGG,
\NRemovedGBM{} GBM), so complete-case analysis does not distort the class balance.

\textbf{No target leakage.} Leakage occurs when information that would not be available at
prediction time, or that encodes the target, enters the model \cite{kaufman2012}. Two columns do
exactly that: the TCGA project coincides with the grade for all \NProjectGrade{} patients, and the
histological diagnosis determines it (every glioblastoma is GBM and every other diagnosis is LGG).
Both are excluded from all models. The EDA (Section~\ref{sec:prdl-eda}) quantifies the effect of
ignoring this rule.

\textbf{Correct variable types.} Race is a nominal variable whose integer codes have no order, so
it is one-hot encoded; treating the codes as a number would impose an arbitrary linear effect.
Mutations are binary and are kept as 0/1 indicators. All transformations that learn from data
(standardisation, encoding) are fitted inside the cross-validation (Section~\ref{sec:prdl-setup}).

% =====================================================================
% 3. EXPLORATORY DATA ANALYSIS (EDA)
% =====================================================================
\section{Exploratory Data Analysis (EDA)}\label{sec:prdl-eda}

\subsection{Introduction}

The EDA has a concrete purpose in a supervised problem: to find which variables separate the
classes or relate to the target, to detect problems (leakage, redundancy, rare categories) before
modelling, and to set realistic expectations. Its structure follows the course EDA notebook. Because
almost every predictor here is binary, the techniques were adapted, and an unsupervised analysis of
the mutation profiles was added. All EDA statistics describe the full dataset; none of them is used
to select predictors, because selecting variables with the labels outside the validation loop biases
the error estimate \cite{ambroise2002}.

\subsection{Techniques}

\textbf{Tests and effect sizes.} Age is not normally distributed within either class (Shapiro--Wilk
test), so the two classes are compared with the Mann--Whitney $U$ test \cite{mann1947} rather than
a $t$-test. Its statistic has a direct link with classification: $U/(n_1 n_0)$ is the probability
that a random GBM patient is older than a random LGG patient, which is the area under the ROC curve
(AUC) of age used alone \cite{hanley1982}. Binary predictors are compared with Fisher's exact test
\cite{fisher1922}, which, unlike the $\chi^2$ test, remains valid when some cells of the $2\times2$
table are small, as happens for rarely mutated genes. Since a $p$-value does not measure the size of
an effect, every test is reported with an effect size (Cohen's $d$, odds ratio, single-feature AUC
with a bootstrap interval \cite{efron1993}). With \NTests{} simultaneous tests, $p$-values are
corrected with Bonferroni (control of any false positive) and Benjamini--Hochberg (control of the
expected proportion of false discoveries, less conservative) \cite{benjamini1995}.

\textbf{Redundancy.} For binary variables the Pearson correlation is the phi coefficient. The
variance inflation factor, $\mathrm{VIF}_j = 1/(1-R_j^2)$, measures how well a predictor is
explained by all the others; values above 5--10 signal problematic collinearity \cite{james2021}.

\textbf{Unsupervised techniques.} PCA finds orthogonal directions of maximum variance after
centring the data; standardising 0/1 variables gives a rare mutation the same weight as a frequent
one, so both versions are compared. Euclidean geometry treats two patients as similar when they
share the \emph{absence} of mutations, which dominates sparse data. The Jaccard distance,
$d_J(A,B) = 1 - |A\cap B|/|A\cup B|$, ignores shared absences; it is used for metric MDS, whose
fit is measured with Kruskal's stress \cite{kruskal1964}, and for t-SNE, a non-linear map that
preserves local neighbourhoods \cite{vandermaaten2008}. $k$-means minimises the within-cluster
variation; on 0/1 data the squared Euclidean distance equals the number of genes in which two
patients differ. Hierarchical clustering needs a dissimilarity and a linkage; Ward's linkage merges
the groups that least increase the within-cluster variance \cite{ward1963}. Partitions are compared
with the silhouette \cite{rousseeuw1987} and the adjusted Rand index (ARI), which is 0 for chance
agreement and 1 for identical partitions \cite{hubert1985}.

\subsection{Results \& Conclusions}

\textbf{The classes differ strongly in a few variables.} GBM patients are older (median
\AgeMedianGBM{} vs \AgeMedianLGG{} years; Cohen's $d = \AgeCohenD$; AUC \AUCAge{} \AUCAgeCI), and the
largest single difference is IDH1, mutated in \PrevIDHOneLGG\% of LGG and \PrevIDHOneGBM\% of GBM
patients (AUC \AUCIDHOne{} \AUCIDHOneCI; Figure~\ref{fig:prdl-prevalence}). PTEN and EGFR mutations
are more frequent in GBM, and CIC and NOTCH1 almost exclusive to LGG (NOTCH1:
\PrevNOTCHLGG\% vs \PrevNOTCHGBM\%), which anticipates the separation problems of logistic
regression (Section~\ref{sec:prdl-results}). After correction, \NSigBonf{} of \NTests{} differences
remain significant with Bonferroni and \NSigBH{} with Benjamini--Hochberg; most of the remaining
genes have AUCs close to 0.5, so statistical significance does not imply usefulness.

\begin{figure}[H]
\centering
\includegraphics[width=0.62\textwidth]{eda_mutation_prevalence.png}
\caption{Percentage of patients with each mutation in LGG and GBM (sorted by the difference).}
\label{fig:prdl-prevalence}
\end{figure}

\textbf{The redundancy is moderate and biologically structured.} The strongest correlations are
age--IDH1 ($r = \CorrAgeIDH$), TP53--ATRX ($r = \CorrTPATRX$) and CIC--FUBP1 ($r = \CorrCICFUBP$),
and the maximum VIF is \MaxVIFEDA, below the usual thresholds. Correlated mutations do not make linear
models unstable here, but they violate the independence assumption of naive Bayes.

\textbf{Leakage would make the problem look solved.} A quick cross-validated logistic regression
reaches an AUC of \BaseAUCAll{} with the 23 legitimate predictors and \BaseAUCAgeIDH{} with only age
and IDH1, but \BaseAUCLeak{} when the histological diagnosis is added. The last number measures
the leakage, not the information in the data.

\textbf{The unsupervised structure is the IDH axis.} PCA does not compress the profiles: \NCompCentred{}
centred (or \NCompScaled{} scaled) components are needed to explain 80\% of the variance, and PC1,
which explains \PCOneCentred\% of the centred variance, is the contrast between IDH1/TP53/ATRX and
PTEN/EGFR. With the Jaccard distance, \PctJaccardOne\% of the patient pairs share no mutation at all
($d_J = 1$), and the MDS map distorts the distances considerably (stress \MDSStress). Hierarchical
clustering with Jaccard distance fails: the \NNoMutation{} patients without any mutation are at
distance 1 from everybody, so complete linkage ties at 1 and merges all patients into one cluster,
while average and single linkage only isolate those \NNoMutation{} patients. This is a property of the
distance, not a bug, and it illustrates why the dissimilarity must be chosen for the data.
$k$-means and Ward's method on the 0/1 profiles agree well. The number of clusters was fixed before
looking at the result by a rule (highest agreement between the two algorithms), which gave
$k = \KMain$ (ARI \ARIKMWardMain). The two clusters are essentially IDH-wildtype
(\CluMainAN{} patients, \CluMainAGBM\% GBM, median age \CluMainAAge) and IDH-mutant (\CluMainBN{} patients,
\CluMainBGBM\% GBM, median age \CluMainBAge), with ARI \ARIGradeMain{} against the grade. As an
exploratory, secondary solution, $k = \KSecond$ splits the IDH-mutant group into a TP53/ATRX profile
and a CIC/FUBP1 profile with \CluSecondCGBM\% GBM (Figure~\ref{fig:prdl-clusters}), a structure that
mirrors the TCGA description of lower-grade gliomas: IDH-mutant tumours without 1p/19q codeletion
nearly all had TP53 mutations (94\%) and ATRX inactivation (86\%), while codeleted ones carried CIC,
FUBP1 and NOTCH1 mutations \cite{tcga2015lgg}. The clusters thus recover, from point mutations only,
the astrocytoma/oligodendroglioma distinction that the classification makes with the codeletion. Because the
cluster structure was not imposed by the grade, its agreement with the grade supports H3, and its
lack of perfect agreement is consistent with the outdated labels discussed in Section 2.1.

\begin{figure}[H]
\centering
\includegraphics[width=0.85\textwidth]{unsup_cluster_profiles.png}
\caption{Mutation prevalence in each $k$-means cluster for the main ($k=\KMain$) and the secondary
($k=\KSecond$) solution. The grade was not used to build the clusters.}
\label{fig:prdl-clusters}
\end{figure}

\textbf{Conclusions for modelling.} (i) Accuracy alone is misleading because the majority class
already gives \MajorityAcc; AUC, sensitivity and specificity are reported. (ii) Age and IDH1 carry most
of the signal, so large gains from complex models are not expected. (iii) Rare categories (one patient
of one race, mutations present in a single class) require penalised models and encoders that tolerate
unseen categories. (iv) The diagnosis and the project must never be predictors.

% =====================================================================
% 4. EXPERIMENTAL SETUP
% =====================================================================
\section{Experimental setup}\label{sec:prdl-setup}

\subsection{Introduction}

With \NClean{} patients, setting aside a single test set would waste data and make the estimate
depend on one random split. The alternative used here is \textbf{nested cross-validation}, the scheme
of the course example: an outer loop of 10 folds estimates the performance, and inside each outer
training part an inner loop of 3 folds (\texttt{GridSearchCV}) chooses the hyperparameters. The tuned
model is refitted on the whole outer training part and predicts the held-out fold, which played no
role in any decision. This matters because the inner cross-validation score of the selected
configuration is optimistically biased: it is the maximum over many noisy estimates. Varma and Simon
showed that ``a nested CV procedure provides an almost unbiased estimate of the true error''
\cite{varma2006}, and Cawley and Talbot that ``model selection must be treated as an integral part of
the model fitting process and performed afresh every time a model is fitted'' \cite{cawley2010}.
The choice of $k = 10$ outer folds follows the usual bias--variance compromise of $k$-fold
cross-validation \cite{james2021}.

Four further decisions complete the protocol:
\begin{itemize}
  \item \textbf{Everything that learns from data is inside the loop.} Standardisation and one-hot
        encoding are steps of a pipeline that is refitted on each training part, so the statistics of
        a test fold never influence its own preprocessing.
  \item \textbf{Stratified folds for classification}, so that each fold keeps the
        \PctLGG\%/\PctGBM\% class proportions; plain shuffled folds for regression.
  \item \textbf{Metrics with a reference point.} Classification: ROC-AUC (threshold-free; also the
        inner selection criterion), and at the default threshold 0.5, sensitivity (share of GBM
        detected), specificity (share of LGG correctly identified), balanced accuracy and accuracy,
        all compared with a model that predicts the class proportions. Regression: MAE and RMSE (in
        years) and $R^2$, compared with predicting the mean age; RMSE is the inner criterion because
        it matches the squared-error loss minimised by most of the models.
  \item \textbf{Uncertainty and comparison rule.} Every patient receives one out-of-fold prediction.
        Metrics are computed on these predictions with 95\% percentile bootstrap intervals (2000
        resamples of patients, the same for every model) \cite{efron1993}. A model is \emph{not}
        declared better than another when their intervals overlap; because two models are evaluated
        on the same patients, a paired bootstrap of the difference is reported as a more sensitive
        complement.
\end{itemize}

Nested cross-validation estimates how well the \emph{procedure} (model family plus tuning)
generalises, not the performance of one fitted model. The models interpreted in
Section~\ref{sec:prdl-results} are therefore obtained by applying the same procedure once to all
patients. Even with these precautions, the evaluation is internal: the reporting guidance TRIPOD+AI
requires evaluation data ``distinct from the data used to train the model, tune hyperparameters, or
do model selection'' \cite{collins2024}, which nested cross-validation satisfies within the cohort,
but it cannot replace validation on an external cohort.

% =====================================================================
% 5. MACHINE LEARNING MODELS
% =====================================================================
\section{Machine Learning models}

\subsection{Model selection}

The models were chosen to span the course from the most constrained to the most flexible, so that
the comparison answers a question (how much flexibility does this problem need?) rather than
searching for the highest score. For each family, the hyperparameters searched are those that control
its bias--variance trade-off \cite{james2021,hastie2009}.

\textbf{Linear and penalised linear models.} Multiple linear regression estimates
$y = \beta_0 + \sum_j \beta_j x_j + \varepsilon$ by least squares; with 0/1 predictors each
$\beta_j$ is the difference in mean age between patients with and without mutation $j$, holding the
others fixed. Its inference assumes independent, homoscedastic, normal errors, which are checked with
residual diagnostics. Ridge \cite{hoerl1970} and Lasso \cite{tibshirani1996} add the penalties
$\lambda\sum_j\beta_j^2$ and $\lambda\sum_j|\beta_j|$; Ridge shrinks all coefficients, while the
corners of the L1 constraint set some exactly to zero, i.e.\ Lasso selects variables. The penalty
$\lambda$ (\texttt{alpha}) is tuned, and inputs are standardised so that the penalty treats all
coefficients alike. Logistic regression models the log-odds of GBM as a linear function, with an L1
or L2 penalty and $C = 1/\lambda$ tuned. The penalty is also needed for a second reason: when a
predictor occurs in only one class, maximum-likelihood estimates do not exist (complete or
quasi-complete separation) \cite{albert1984}; NOTCH1 is such a predictor here.

\textbf{Naive Bayes} applies Bayes' rule assuming the predictors are independent given the class.
Age is discretised into quantile bins so that every column is Bernoulli; the number of bins and the
Laplace smoothing are tuned. The independence assumption is false for correlated mutations, which
often leaves the ranking of patients (and hence the AUC) reasonable while making the probabilities
over-confident \cite{hand2001}.

\textbf{Trees and ensembles.} A tree partitions the predictor space into boxes by recursive binary
splits chosen to reduce impurity (Gini index or entropy/deviance for classification, squared error for
regression) \cite{breiman1984}. A large tree overfits, so it must be pruned. The regression tree is
limited in advance by maximum depth and minimum leaf size; the classification tree is grown and then
pruned by cost-complexity pruning, which minimises impurity $+\ \alpha\cdot$(number of leaves) with
$\alpha$ (\texttt{ccp\_alpha}) tuned. Random forests average many trees grown on bootstrap samples,
trying a random subset of predictors (\texttt{max\_features}) at each split to decorrelate them
\cite{breiman2001}; averaging reduces variance, so the trees need no pruning, and the minimum leaf size
is tuned. Gradient boosting fits small trees sequentially to the gradient of the loss of the current
ensemble; XGBoost \cite{chen2016} is the implementation used, with the number of trees, the shrinkage
(learning rate) and the depth tuned.

\textbf{Support vector machines} find the maximum-margin hyperplane allowing violations penalised by
$C$, and the kernel trick gives non-linear boundaries; linear and RBF kernels are compared, with $C$
and the RBF reach $\gamma$ tuned \cite{cortes1995}. Probabilities for the AUC are obtained by Platt
scaling \cite{platt1999}. Support vector regression fits a function within an $\varepsilon$-tube,
ignoring errors smaller than $\varepsilon$ years.

\textbf{Baselines.} The mean age and the class proportions give the reference that any useful model
must beat.

\subsection{Complementary analyses}

\textbf{Multiclass classification.} Logistic regression extends to $K$ classes by giving every class
$k$ its own linear score $z_k = \beta_{k0} + \boldsymbol\beta_k^\top\mathbf x$ and turning the scores
into probabilities with the softmax function, $P(Y=k\mid\mathbf x) = e^{z_k}/\sum_{j} e^{z_j}$
(multinomial regression) \cite{james2021}. Trees and random forests handle several classes directly,
with the Gini index or the entropy computed over the $K$ class proportions of each node, and XGBoost
fits one tree per class at every boosting round on the softmax loss \cite{chen2016}. The main metric is
the multiclass AUC of Hand and Till, the mean over all pairs of classes $(i,j)$ of the AUC of
$P(Y=i)$ separating class $i$ from class $j$, averaged with the reverse comparison; like the binary AUC
it needs no threshold and does not depend on the class proportions \cite{handtill2001}. Balanced
accuracy (mean recall over the classes; 0.25 for a constant prediction with four classes) and the macro
F1 score describe the predicted class.

\textbf{SHAP values.} Permutation importance says how much performance is lost without a predictor,
but not in which direction a predictor pushes each prediction. SHAP values \cite{lundberg2017} split one
prediction $f(\mathbf x)$ into additive contributions, $f(\mathbf x) = \phi_0 + \sum_j \phi_j$, where
$\phi_0$ is the average prediction and $\phi_j$ is the Shapley value of predictor $j$
\cite{shapley1953}: the average change in the prediction when $j$ is added to every possible subset of
the other predictors. For trees and tree ensembles the values are computed exactly by TreeSHAP
\cite{lundberg2020}; for a linear model with independent predictors they reduce to
$\phi_j = \beta_j(x_j - \bar x_j)$. SHAP values describe the fitted model, not the disease: correlated
predictors share the credit in a way that depends on the algorithm, and a large value is not evidence of
a causal effect.

\textbf{Causal analysis.} A prediction model may use any variable that carries information; a causal
question requires deciding \emph{which} variables to adjust for, and that decision comes from a causal
diagram (a directed acyclic graph, DAG) rather than from the fit of the model
\cite{pearl2009,huntingtonklein2021}. On a path from an exposure to an outcome, a \emph{confounder} is a
common cause and must be adjusted for; a \emph{mediator} lies on the causal path, and adjusting for it
removes the part of the effect that goes through it; a \emph{collider} is a common effect, and adjusting
for it creates a spurious association. With an exposure $E$, a mediator $M$, an outcome $Y$ and the same
covariates, three least-squares regressions give the total effect $c$ ($Y$ on $E$), the effect $a$ of
$E$ on $M$, and the direct effect $c'$ and the effect $b$ of $M$ ($Y$ on $E$ and $M$); for ordinary
least squares on the same patients the total effect splits exactly as $c = c' + ab$, where $ab$ is the
indirect effect through the mediator \cite{baron1986,mackinnon1995}. The split is causal only under the
assumptions of the diagram, in particular that no unmeasured variable affects both the mediator and the
outcome.

% =====================================================================
% 6. RESULTS
% =====================================================================
\section{Results}\label{sec:prdl-results}

\subsection{Evaluation methodology}

All numbers below are out-of-fold estimates from the nested cross-validation of
Section~\ref{sec:prdl-setup}, with 95\% bootstrap intervals in brackets. The coefficients, trees and
importances that are interpreted come from the final models fitted on all patients with the same
tuning procedure; their in-sample fit is not used as evidence of performance.

\subsection{Discussion}

\textbf{Regression: the profile explains part of the variance of age, and every model is
equivalent.} Predicting the mean age gives a mean absolute error of \RegBaseMAE{} years. All models
reduce it to between \RegMinMAE{} and \RegMaxMAE{} years and explain between \RegMinRTwo{} and
\RegMaxRTwo{} of the variance (Table~\ref{tab:prdl-reg}). The paired bootstrap confirms that every
model improves on the baseline. The \RegBest{} has the best point estimates (MAE \RegBestMAE{}
\RegBestMAECI{} years, $R^2$ \RegBestRTwo{} \RegBestRTwoCI), but its interval overlaps those of the
other \RegNTies{} models, and even the paired difference with linear regression includes zero
(\RegDiffOLSBest{} \RegDiffOLSBestCI{} years). By the rule fixed in advance, no model is declared
better: the additional flexibility of ensembles and kernels buys nothing that a linear model cannot
capture, which is expected when the predictors are binary and their effects nearly additive.

\begin{table}[H]
\centering
\caption{Regression of age at diagnosis: out-of-fold metrics with 95\% bootstrap intervals, and
$R^2$ per outer fold (mean $\pm$ SD).}
\label{tab:prdl-reg}
\resizebox{\textwidth}{!}{\input{tables/reg_model_comparison.tex}}
\end{table}

The interpretation is more interesting than the ranking. The multiple linear regression attributes
the largest effects to IDH1 (\OLSIDHOne{} years, 95\% CI \OLSIDHOneCI), IDH2 (\OLSIDHTwo{} years,
\OLSIDHTwoCI) and ATRX (\OLSATRX{} years); only \OLSNSigBH{} coefficients remain significant after
Benjamini--Hochberg correction, the maximum VIF is \OLSMaxVIF{}, and no patient has a Cook's distance
close to 1. The residual tests reject exact normality and homoscedasticity (Shapiro $p = \OLSShapirop$,
Breusch--Pagan $p = \OLSBPp$), so the coefficient intervals are approximate. One coefficient shows a
limitation of the data: the only American Indian or Alaska Native patient has leverage 1, so the model
fits that patient exactly and the coefficient (\OLSNative{} years, \OLSNativeCI) is meaningless.
Lasso, with the penalty chosen by cross-validation, sets \LassoNZero{} of \LassoNCoef{} coefficients
to zero and keeps the same leading predictors.

The out-of-fold predictions explain why the models agree (Figure~\ref{fig:prdl-oof}): they cluster
around two values. Their correlation with the grade is \RegCorrGrade, but with the IDH status (IDH1 or
IDH2 mutated) it is stronger, \RegCorrIDH{} (IDH-mutant patients are predicted younger), and Table~\ref{tab:prdl-reg-idh} separates the two: within GBM, the
\RegNGBMIDHmut{} IDH-mutant patients are predicted at \RegPredGBMIDHmut{} years on average, against
\RegPredGBMIDHwt{} for IDH-wildtype GBM, and within LGG the split is \RegPredLGGIDHmut{} vs
\RegPredLGGIDHwt{} years. The models predict the age from the IDH status, not from the grade; the
predictions correlate with the grade only because the grade is associated with IDH. A first draft of
this report read the two groups as the grade (Appendix~A); the causal analysis below quantifies how
much of the IDH--age association actually goes through the grade.

\begin{table}[H]
\centering
\caption{Mean observed and out-of-fold predicted age (years) of the \RegBest{} by grade and IDH status
(neither was an input of the model).}
\label{tab:prdl-reg-idh}
\small
\input{tables/reg_pred_by_grade_idh.tex}
\end{table}

\begin{figure}[H]
\centering
\includegraphics[width=0.9\textwidth]{reg_oof_predictions.png}
\caption{Out-of-fold predictions of the \RegBest{}: observed vs predicted age and residuals, coloured by
grade (not an input of the model).}
\label{fig:prdl-oof}
\end{figure}

\textbf{Classification: all models reach similar AUCs; simple models are as good as
ensembles.} Against a baseline with AUC \ClfAUCBase{} and accuracy \ClfBaseAcc, the models reach
AUCs between \ClfMinAUC{} and \ClfMaxAUC{} (Table~\ref{tab:prdl-clf}, Figure~\ref{fig:prdl-roc}). The
\ClfBest{} has the highest AUC (\ClfBestAUC{} \ClfBestAUCCI), with sensitivity \ClfBestSens{} and
specificity \ClfBestSpec{} at the default threshold, but the intervals of all \ClfNTies{} other models
overlap it, so by the pre-specified rule there is no winner. The paired comparison is more
informative: logistic regression (\ClfDiffLR) and XGBoost (\ClfDiffXGB) are indistinguishable from the
random forest, while naive Bayes (\ClfDiffNB), the SVM (\ClfDiffSVM) and the pruned tree
(\ClfDiffTree) are slightly but consistently worse. The two rules answer different questions, and
both are reported: overlapping marginal intervals say that the differences are small relative to the
uncertainty of each estimate; the paired test says that some of these small differences are
systematic. Either way, H1 is supported: a penalised logistic regression on these 23 predictors is as
good as any ensemble.

\begin{table}[H]
\centering
\caption{Classification of LGG vs GBM: out-of-fold metrics with 95\% bootstrap intervals (threshold
0.5 for sensitivity, specificity and balanced accuracy).}
\label{tab:prdl-clf}
\resizebox{\textwidth}{!}{\input{tables/clf_model_comparison.tex}}
\end{table}

\begin{figure}[H]
\centering
\includegraphics[width=0.55\textwidth]{clf_roc_curves.png}
\caption{Out-of-fold ROC curves of the six classifiers.}
\label{fig:prdl-roc}
\end{figure}

Three results deserve a critical reading.
\emph{(i) What the models learn.} Permutation importance ranks IDH1 far above every other predictor
(mean AUC drop \ClfPermIDH{}), followed by age; the pruned tree (depth \ClfTreeDepth, \ClfTreeLeaves{}
leaves) first splits on IDH1 and then on IDH2, TP53 and age. The models reproduce the IDH axis found
without labels in the EDA, which is consistent with the role of IDH in the WHO classification. The
\ClfFN{} GBM patients predicted as LGG and the \ClfFP{} LGG patients predicted as GBM by the
\ClfBest{} cannot be explained from these data alone; part of them may be label noise from the
pre-2021 diagnoses (Section 2.1).
\emph{(ii) Separation and the meaning of coefficients.} The final logistic regression chose an
almost unpenalised L1 model ($C = \LRFinalC$) although the outer folds chose smaller values: the
inner-CV surface is flat (\LRFlatWithin{} of \LRFlatTotal{} configurations within 0.005 AUC of the
best), so the choice among them is close to arbitrary. With a weak penalty, NOTCH1, mutated only in
LGG patients (in TCGA, a mutation of 1p/19q-codeleted gliomas \cite{tcga2015lgg}), receives a coefficient (\LRCoefNOTCH{} log-odds per SD) comparable to that of IDH1
(\LRCoefIDH): a direct consequence of quasi-separation \cite{albert1984}, not evidence that NOTCH1 is
as important as IDH1. The odds ratio of age, \LRORAge{} per standard deviation, is more stable.
\emph{(iii) Ranking vs calibration.} Naive Bayes reaches a competitive AUC, but \NBExtreme\% of its
out-of-fold probabilities are below 0.02 or above 0.98, against \LRExtreme\% for logistic regression.
Because it counts the evidence of correlated mutations (IDH1, ATRX, TP53) several times, its
probabilities are over-confident even when its ranking is good \cite{hand2001}. If the probabilities
were to be used, for example to choose a clinical threshold, this difference would matter more than
the AUC.

\textbf{Comparison with previous work.} On the same 839 TCGA patients, \cite{tasci2022} reported an
accuracy of 87.606\% with a voting ensemble and feature selection, and \cite{sanchezmarques2024}
0.882 with oversampling. The accuracies obtained here (between \ClfMinAcc{} and \ClfMaxAcc, MLLB report) are of the same order with standard models and no resampling. A strict comparison is not
possible because the validation schemes differ. In \cite{sanchezmarques2024} the hyperparameters were
``fine-tuned [...] using an 80-20 stratified holdout setting'' and performance was then evaluated with
10-fold cross-validation; if the tuning split is drawn from the same patients, the tuned values have
seen the evaluation folds, and the estimate is optimistic \cite{varma2006,cawley2010}.
In \cite{tasci2022}, instead, no tuning was done: ``the default values were assigned as the
corresponding parameter values for the utilized classifiers'', so their estimate is not inflated by
hyperparameter selection, but their classifiers were not tuned either. The nested cross-validation
used here tunes every model while keeping the evaluation folds out of the tuning, at the cost of
slightly lower but honest numbers.

\textbf{Multiclass: the mutations recognise the families, except the mixed gliomas.} Against a baseline
macro AUC of \McBaseAUC, the four models reach between \McMinAUC{} and \McMaxAUC{}
(Table~\ref{tab:prdl-mc}). The \McBest{} has the highest value (\McBestAUC; balanced accuracy
\McBestBAcc, macro F1 \McBestFOne), the intervals of the \McNTies{} other models overlap it, and in the
paired comparison only the \McWorse{} is worse. The probability of glioblastoma ranks the grade almost as
well as the binary models (AUC \McGradeAUC), so the three extra classes cost nothing for the grade. The
confusion matrix (Figure~\ref{fig:prdl-mc}) shows where the information lies:
\begin{itemize}
  \item \textbf{Glioblastoma} is recognised best (recall \McRecallGBM, one-vs-rest AUC \McAUCGBM),
        followed by \textbf{oligodendroglioma} (\McRecallOligo, \McAUCOligo) and \textbf{astrocytoma}
        (\McRecallAstro, \McAUCAstro).
  \item \McAstroAsGBM{} astrocytomas are predicted as glioblastoma; \McAstroGBMwt{} of them are
        IDH-wildtype, and they are all the IDH-wildtype astrocytomas (\McAstroWtTotal) of the cohort.
        In the 2021 classification, for an IDH-wildtype diffuse astrocytic glioma the ``presence of 1 or
        more of 3 genetic parameters (TERT promoter mutation, EGFR gene amplification, combined gain of
        entire chromosome 7 and loss of entire chromosome 10 [+7/$-$10]) appears sufficient to assign''
        the diagnosis of glioblastoma, IDH-wildtype \cite{louis2021}. These markers are not in the data,
        so how many of the \McAstroGBMwt{} patients would be reclassified cannot be checked, but these
        ``errors'' are at least partly the outdated labels of Section~2.1.
  \item No patient is predicted as \textbf{mixed glioma} (recall \McRecallMixed): \McMixedAsAstro\% are
        predicted as astrocytoma, \McMixedAsOligo\% as oligodendroglioma and \McMixedAsGBM\% as
        glioblastoma, and their mean predicted probability is split between the two pure families
        (Figure~\ref{fig:prdl-mc}, right). Their mutation profile is that of one of the two pure
        families. This is what the 2016 WHO classification anticipated: ``the diagnosis of
        oligoastrocytoma is strongly discouraged. Nearly all tumors with histological features suggesting
        both an astrocytic and an oligodendroglial component can be classified as either astrocytoma or
        oligodendroglioma using genetic testing'' \cite{louis2016}. Here, point mutations alone already
        push most of them towards one of the two families.
\end{itemize}
The unsupervised analysis of Section~\ref{sec:prdl-eda} anticipated this structure without labels: the
CIC/FUBP1 cluster is \McCThreeOligo\% oligodendrogliomas but also contains \McCThreeMixed{} mixed
gliomas, and the TP53/ATRX cluster is \McCTwoAstro\% astrocytomas with \McCTwoMixed{} mixed gliomas.

\begin{table}[H]
\centering
\caption{Multiclass classification of the histological family: out-of-fold metrics with 95\% bootstrap
intervals (\McNAstro{} astrocytomas, \McNMixed{} mixed gliomas, \McNOligo{} oligodendrogliomas, \McNGBM{}
glioblastomas).}
\label{tab:prdl-mc}
\resizebox{\textwidth}{!}{\input{tables/mc_model_comparison.tex}}
\end{table}

\begin{figure}[H]
\centering
\includegraphics[width=\textwidth]{mc_confusion.png}
\caption{\McBest: out-of-fold confusion matrix (colour: share of each true family) and mean predicted
probability of each family by true family.}
\label{fig:prdl-mc}
\end{figure}

\textbf{What the winning models use: SHAP.} SHAP values were computed for the three binary models
that the paired comparison could not separate, and for the multiclass \ShapMcModel{}
(Figures~\ref{fig:prdl-shap} and~\ref{fig:prdl-shap-mc}). Their global ranking agrees with the
permutation importance (Spearman $\rho$ = \ShapRhoRF{} for the random forest, \ShapRhoLR{} for logistic
regression and \ShapRhoXGB{} for XGBoost), and IDH1 takes \ShapIDHRF\% of the total mean $|$SHAP$|$ of
the random forest and \ShapIDHXGB\% of XGBoost. SHAP adds three things that the permutation importance
could not show:
\begin{itemize}
  \item \textbf{The direction and size of each effect, patient by patient.} NOTCH1 is the second input of
        logistic regression (\ShapNOTCHLR\% of its total) but negligible in the random forest
        (\ShapNOTCHRF\%): the few NOTCH1-mutated patients receive a very large contribution towards LGG,
        the quasi-separation effect discussed above made visible.
  \item \textbf{The shape of the effect of age.} In logistic regression it is a straight line by
        construction; the tree ensembles learn a curve that is almost flat for young patients and rises
        steeply in middle age, changing sign at about \ShapAgeRF{} years in the random forest and
        \ShapAgeXGB{} in XGBoost.
  \item \textbf{What separates the families.} In the multiclass forest, CIC weighs most for
        oligodendroglioma (\ShapCICOligo\% of that class's total, against \ShapCICAstro\% for
        astrocytoma), and ATRX and age more for astrocytoma (\ShapATRXAstro\% and \ShapAgeAstro\%, against
        \ShapATRXOligo\% and \ShapAgeOligo\%), the molecular split described for the TCGA lower-grade
        gliomas \cite{tcga2015lgg}.
\end{itemize}
The ranking agrees with the previous explainability analysis of these data among the sources reviewed:
on the same patients, Sánchez-Marqués et al.\ combined four feature-ranking methods, permutation
importance and SHAP values, and report that the ``IDH1 biomarker was the most important attribute
contributing to the target variable (i.e., glioma grade), regardless of the model used. The second
most important feature was Age at diagnosis'' \cite{sanchezmarques2024}. The same order appears here in
the random forest and XGBoost; logistic regression is the exception, for the quasi-separation reason
above. Tasci et al.\ used feature selection but do not report a ranking of the features
\cite{tasci2022}.

\begin{figure}[H]
\centering
\includegraphics[width=\textwidth]{shap_binary_summary.png}
\caption{SHAP summary of the three tied binary models: one dot per patient, contribution towards GBM
(colour: value of the input, grey = 0 or young, violet = 1 or old). Units: probability for the random
forest, log-odds for logistic regression and XGBoost.}
\label{fig:prdl-shap}
\end{figure}

\begin{figure}[H]
\centering
\includegraphics[width=\textwidth]{shap_multiclass.png}
\caption{Multiclass \ShapMcModel: share of mean $|$SHAP$|$ of each input per family, and contributions
towards astrocytoma and oligodendroglioma.}
\label{fig:prdl-shap-mc}
\end{figure}

\subsection{Causal analysis}\label{sec:prdl-causal}

Figure~\ref{fig:prdl-dag} shows the diagram assumed for Q5. IDH-mutant patients are younger within both
grades (LGG: \CsAgeLGGmut{} vs \CsAgeLGGwt{} years; GBM: \CsAgeGBMmut{} vs \CsAgeGBMwt{} years, with
only \CsNGBMmut{} IDH-mutant GBM), which already suggests that the grade does not carry the whole
association. Adjusted for sex and race, the total effect of an IDH mutation on the age at diagnosis is
\CsTotal{} years (OLS with HC3 standard errors \cite{hayes2007}; bootstrap intervals). Adjusting also
for the grade leaves a direct effect of \CsDirect{} years; the indirect effect through the grade is
\CsIndirect{} years, \CsPropPct\% of the total (Table~\ref{tab:prdl-med}). The identity
$c = c' + ab$ holds exactly, so the result rests on the assumptions, not on the arithmetic.

\begin{figure}[H]
\centering
\begin{minipage}[t]{0.47\textwidth}\centering
\includegraphics[width=\textwidth]{causal_dag.png}
\caption{Assumed causal diagram (dashed: unmeasured).}
\label{fig:prdl-dag}
\end{minipage}\hfill
\begin{minipage}[t]{0.5\textwidth}\centering
\includegraphics[width=\textwidth]{causal_adjustment_sets.png}
\caption{Coefficient of IDH on age for different adjustment sets.}
\label{fig:prdl-adj}
\end{minipage}
\end{figure}

\begin{table}[H]
\centering
\caption{Mediation decomposition of the IDH--age association (adjusted for sex and race; 95\% bootstrap
intervals, 2000 resamples).}
\label{tab:prdl-med}
\small
\input{tables/causal_mediation.tex}
\end{table}

Three readings follow:
\begin{itemize}
  \item \textbf{It corrects H2.} Under the diagram, most of the association between IDH and age does
        not go through the binary grade, in agreement with the regression predictions above.
  \item \textbf{What is adjusted for changes the answer} (Figure~\ref{fig:prdl-adj}). Adjusting for the
        histological family instead of the grade gives \CsFamily{} years, and adding the 18 other
        mutations \CsGenes{} years. The family is itself a consequence of IDH, so adjusting for it is
        over-adjustment for the total effect; whether the other mutations are confounders or mediators
        the diagram does not settle. Adding variables is not ``more control'': the diagram decides
        \cite{huntingtonklein2021,pearl2009}. For the prediction models above the same question does
        not arise, because there any variable that carries information is useful.
  \item \textbf{The direct effect is the fragile part.} If an unmeasured variable affects both the grade
        and the age at diagnosis (U in Figure~\ref{fig:prdl-dag}, for example the copy-number markers
        missing from the data), adjusting for the grade opens the path IDH $\to$ grade $\leftarrow$ U
        $\to$ age, and $c'$ is biased; the total effect does not need that assumption. The decomposition
        also assumes that the effect of IDH is the same in both grades, but the IDH $\times$ grade
        interaction is \CsInter{} years, so it is probably larger in GBM, although that estimate rests
        on \CsNGBMmut{} patients.
\end{itemize}

% =====================================================================
% 7. CONCLUSIONS
% =====================================================================
\section{Conclusions}

\subsection{Conclusions related to the study}

\begin{itemize}
  \item \textbf{H1 (classification) is supported.} Age and 20 point mutations separate LGG from GBM
        with out-of-fold AUCs between \ClfMinAUC{} and \ClfMaxAUC{} under nested cross-validation. Penalised logistic
        regression and XGBoost are statistically indistinguishable from the best random forest;
        naive Bayes, the SVM and the pruned tree are slightly worse in paired comparisons.
  \item \textbf{H2 (regression) is supported with a caveat.} Without the grade, the profile
        explains a share of the variance of age between \RegMinRTwo{} and \RegMaxRTwo{} ($R^2$), with all
        models equivalent; but the predictions
        mostly recover the grade through IDH1, IDH2 and ATRX, so the mutations explain age largely
        through the grade.
  \item \textbf{H3 (unsupervised) is supported.} Without labels, $k$-means and Ward's method
        separate an IDH-wildtype and an IDH-mutant group, related to but not identical with the
        grade; a secondary three-cluster solution separates two IDH-mutant profiles. Jaccard-based
        hierarchical clustering fails on these data because of patients without any mutation.
  \item \textbf{The labels limit the study.} The grades come from pre-2021 criteria and several
        markers of the current classification are missing from the data; the cohort is mostly white
        and only internally validated. The models describe this cohort, not current clinical
        practice.
\end{itemize}

\subsection{Conclusions on what you have learned}

The most useful lessons were methodological. First, a large part of good modelling happens before
the models: the histological diagnosis would have produced a perfect but meaningless classifier, and
only checking how the target was defined revealed it. Second, the choice of a distance or a linkage
is a modelling decision with consequences: the Jaccard distance is the natural choice for sparse
binary data, but it made hierarchical clustering degenerate. Third, ``which model is best'' often has
no answer at the available sample size, and saying so with intervals is more honest than ranking
point estimates; how the comparison is made (marginal intervals or paired differences) must be fixed
in advance. Fourth, a model can rank patients well while giving unreliable probabilities or
coefficients (naive Bayes, logistic regression under quasi-separation), so evaluation must match the
intended use. Finally, the AI assistants used for literature search produced wrong or misattributed
sources that were only caught by cross-checking (Appendix B).


% =====================================================================
% REFERENCES
% =====================================================================
\clearpage
\section{References}
\begingroup
% thebibliography opens its own \section*{References}; it is swallowed here so that
% the numbered banner above is the only heading.
\renewcommand{\section}[2]{}
\bibliographystyle{IEEEtran}
\bibliography{references}
\endgroup

\newpage
% =====================================================================
% APPENDIX A: USE OF AI SCIENTIFIC ASSISTANTS
% =====================================================================
\appendixbanner{APPENDIX A: USE OF AI SCIENTIFIC ASSISTANTS}
\addcontentsline{toc}{section}{Appendix A: Use of AI scientific assistants}

\subsection*{Workflow}

Three AI assistants were used for the theoretical part, each with a defined role and a verification
step (the full log is \texttt{docs/\allowbreak{}AI\_\allowbreak{}USAGE\_\allowbreak{}LOG.\allowbreak{}md} in the repository):
\begin{enumerate}
  \item \textbf{Perplexity} searched the literature with structured prompts that asked, for every
        source, the full reference, the DOI, open access and the exact claim it supports.
  \item \textbf{NotebookLM}, restricted to the downloaded PDFs, produced (a) an inventory of every file
        with the bibliographic data as printed in the document, and (b) for each topic of the report,
        short statements with the file, the page and a verbatim quote of at most 30 words.
  \item \textbf{Claude} decided which sources to keep from the inventory, wrote the text using only the
        quote-supported part of each statement, marked every claim still lacking a quote with a
        visible \texttt{[verify]} tag, and drafted the interpretation of the results, which the author
        checked against the notebooks.
\end{enumerate}
The author checks each tagged claim against the PDF before the final version.

\subsection*{Errors detected (candidates for the AI Error Bounty)}

\begin{enumerate}
  \item \textbf{Off-topic sources.} For a request on glioma classification, Perplexity returned a
        cardiac epigenetic reprogramming paper (twice) and a study on lipid metabolism in glioblastoma
        cell lines, neither supporting any requested point. Detected by comparing title and topic with
        the request; removed.
  \item \textbf{Wrong title and authors, repeated by two tools.} Perplexity and then NotebookLM gave
        one of the downloaded files as ``Machine Learning and Statistical Evaluation of Glioma
        Datasets'' by Garbulowski and Komorowski; the extraction step later attributed its content to
        Sánchez-Marqués et al., and the DOAJ record confirms the latter title and authors
        \cite{sanchezmarques2024}. Two independent AI tools agreeing is therefore not a verification.
  \item \textbf{Misattribution.} Parsons et al.\ (2008, \emph{Science}) was presented as a TCGA
        publication by both tools; it was removed from the TCGA sources.
  \item \textbf{Outdated notion of time.} NotebookLM flagged a 2025 publication year as ``a future
        publication year'', although the inventory was made in October 2026.
  \item \textbf{Wrong DOI.} NotebookLM gave 10.1093/neuonc/noad108 for the EANO guideline on molecular
        diagnostic tools; a Crossref check of all DOIs (script \texttt{report/\allowbreak{}check\_\allowbreak{}references.\allowbreak{}py})
        showed that it belongs to an unrelated article on CAR-T cells, and the correct DOI
        (noad100) was confirmed on the EANO website \cite{sahm2023}. The same check corrected the title
        that NotebookLM gave for \cite{varma2006}.
  \item \textbf{One paper, two identities.} Two uploaded files shared the same DOI; the inventory gave
        one of them a different title and garbled authors, so it was discarded as unreliable, while
        the other was correctly identified later \cite{leiriademendonca2025}.
  \item \textbf{Statements broader than their quotes.} Several extracted statements contained numbers
        or claims absent from the supporting quote (for example, glioblastoma incidence and survival
        figures attributed to a file whose quote only gave lower-grade glioma survival). Only the
        quoted part was used. Two claims of this kind were resolved with a second, targeted query
        that asked for the exact quote: an extracted statement that one study had not tuned its
        models inside the cross-validation was supported only by a quote on feature selection, and the
        targeted query showed that the classifiers used default hyperparameters \cite{tasci2022}; a
        count of reclassified cases, presented as a change from the original TCGA labels, compared
        two re-annotation pipelines instead \cite{leiriademendonca2025}.
  \item \textbf{A plausible but wrong interpretation (Claude).} When drafting the regression results,
        the assistant read the two groups of out-of-fold predictions as the grade (``the models recover
        the grade and predict the typical age of each grade''), supported only by a correlation of
        \RegCorrGrade{} with the grade, and wrote the conclusion of H2 accordingly. \emph{Detection}: the
        causal analysis estimated that only \CsPropPct\% of the IDH--age association goes through the
        grade, which contradicted that reading; checking the alternative explanation showed a stronger
        correlation with the IDH status (\RegCorrIDH) and the same split of predictions within each
        grade (Table~\ref{tab:prdl-reg-idh}). \emph{Correction}: the text, the conclusion of H2 and the
        regression notebook (which now reports both correlations) were rewritten. A correlation between
        a prediction and a variable that was not an input does not show that the model uses that
        variable when it is correlated with another one.
\end{enumerate}

\subsection*{AI Error Bounty sheets}

The strongest errors are documented below in the format of the course guideline. They were found in
the verification loop described above, in which every output of the assistants was checked against its
source, the data or a second run; each check is reproduced by the author from the steps in
\texttt{docs/\allowbreak{}BOUNTY\_\allowbreak{}CHECKS.\allowbreak{}md}.

\subsubsection*{Sheet P1. A wrong DOI pointing to another paper (NotebookLM)}
\begin{description}[style=nextline, leftmargin=1.2em]
  \item[Prompt] Inventory of the uploaded PDFs: one row per file with the bibliographic data ``as printed
        in the document'', writing ``not stated'' instead of guessing.
  \item[Faulty response] For the EANO guideline on molecular diagnostic tools (Sahm et al., 2023) the
        DOI 10.1093/neuonc/noad108.
  \item[Proof] The script \texttt{report/\allowbreak{}check\_\allowbreak{}references.\allowbreak{}py} compares the title and year of every DOI
        with the bibliography. The DOI resolved to an unrelated article on CAR-T cells in
        glioblastoma; the EANO page gives 10.1093/neuonc/noad100.
  \item[Correction] DOI replaced \cite{sahm2023}; every DOI of the bibliography is now checked with the
        script.
  \item[Reproduction] \todoverify{author: open both DOIs and record the titles}
  \item[Reflection] A DOI looks like hard data, but here it was generated, not read: the metadata of a
        source need the same verification as its claims.
\end{description}

\subsubsection*{Sheet P2. The same wrong authors from two tools (Perplexity and NotebookLM)}
\begin{description}[style=nextline, leftmargin=1.2em]
  \item[Prompt] Perplexity: machine-learning studies on glioma grading with clinical and mutation data,
        with full reference and DOI. NotebookLM: the inventory prompt of Sheet P1.
  \item[Faulty response] Both tools gave one of the downloaded papers as ``Machine Learning and
        Statistical Evaluation of Glioma Datasets'' by Garbulowski and Komorowski.
  \item[Proof] The claim-extraction step of NotebookLM attributed the content of the same file to
        Sánchez-Marqués et al.; the DOAJ record of its DOI confirms that title and authors. The two outputs of the same tool contradicted each other.
  \item[Correction] Reference corrected \cite{sanchezmarques2024}; agreement between AI tools is no
        longer treated as verification.
  \item[Reproduction] \todoverify{author: read the first page of 21.pdf}
  \item[Reflection] Two assistants agreeing is not two independent sources: the second may simply have
        repeated the first.
\end{description}

\subsubsection*{Sheet P3. Statements broader than their supporting quotes (NotebookLM)}
\begin{description}[style=nextline, leftmargin=1.2em]
  \item[Prompt] For each topic of the report: short statements, each with file, page and a verbatim
        quote of at most 30 words.
  \item[Faulty response] (a) ``Neither paper performed nested CV'' for Tasci et al., supported by a quote
        about feature selection; (b) ``811 samples maintained consistent classification, 134 changed''
        presented as the change from the original TCGA labels to the 2021 classification.
  \item[Proof] Each statement was compared with its own quote, and the quotes did not contain the
        claims. A second, targeted query showed that (a) Tasci et al.\
        used ``default values'' for the hyperparameters \cite{tasci2022}, and (b) the two compared
        classifications were two re-annotation pipelines, neither of them the original label
        \cite{leiriademendonca2025}.
  \item[Correction] Only the quoted content is used in Sections~2 and~6, with the caveats.
  \item[Reproduction] \todoverify{author: read 19.pdf p.~5 and 42.pdf p.~3}
  \item[Reflection] Asking for a quote is useful only if the quote is then read against the statement;
        the summary around a correct quote can still be wrong.
\end{description}

\subsubsection*{Sheet P4. A plausible but wrong interpretation of a figure (Claude)}
\begin{description}[style=nextline, leftmargin=1.2em]
  \item[Prompt] Regression notebook with the out-of-fold predictions and their interpretation; later,
        drafting of the results section of this report.
  \item[Faulty response] ``The models recover the grade through IDH1, IDH2 and ATRX and predict the
        typical age of each grade'', based on a correlation of \RegCorrGrade{} between the predictions
        and the grade; the conclusion of H2 said that the mutations explain age ``largely through the
        grade''.
  \item[Proof] The causal analysis estimated that only \CsPropPct\% of the
        IDH--age association goes through the grade. Testing the alternative explanation:
        the predictions correlate \RegCorrIDH{} with the IDH status, and within GBM the IDH-mutant
        patients are predicted at \RegPredGBMIDHmut{} years against \RegPredGBMIDHwt{} for IDH-wildtype
        (Table~\ref{tab:prdl-reg-idh}).
  \item[Correction] Notebook 04 reports both correlations and the grade $\times$ IDH table; the
        discussion and the conclusion of H2 were rewritten.
  \item[Reproduction] \todoverify{author: run notebook 04, section 6, and compare}
  \item[Reflection] A correlation between a prediction and a variable that was not an input does not
        show that the model uses that variable when another variable is correlated with both.
\end{description}

\subsection*{Reflection}

The verification loop changed how the sources were used: the bibliography records, for every entry,
how its metadata were checked, and no claim enters the text on the authority of an assistant alone.


\newpage

% =====================================================================
% APPENDIX B: ENSEMBLE LEARNING
% =====================================================================
\appendixbanner{APPENDIX B: ENSEMBLE LEARNING}
\addcontentsline{toc}{section}{Appendix B: Ensemble learning}

\subsection*{Classical learning paradigms}

Supervised learning works from examples that already carry the right answer: each
training case pairs a set of inputs with a known target, and the model adjusts itself to
close the gap between its predictions and those targets \cite{hastie2009}. The task is
regression when the target is a number, such as a price or a glucose level, and
classification when it is a category, such as spam versus legitimate email. Typical
algorithms include linear and logistic regression, k-nearest neighbors, decision trees,
support vector machines, and neural networks trained with backpropagation. Applications
include recognizing handwritten digits from pixel data, filtering spam by word frequency,
and predicting heart attack risk from clinical variables \cite{hastie2009}.

Unsupervised learning gets no target at all. Instead of minimizing an error against a
known answer, it looks for structure already sitting in the data: groups of similar
points, directions of high variance, or an underlying distribution \cite{hastie2009}.
Without ground truth to check against, judging an unsupervised model usually comes down
to internal measures like cluster compactness rather than a clean loss function. The two
main families are clustering, grouping similar observations with methods like k-means,
hierarchical clustering, or DBSCAN, and dimensionality reduction, which compresses many
variables into a few while keeping most of the useful information, typically through PCA
or its kernel variants. It is used for grouping gene expression profiles, segmenting
customers for marketing, and flagging anomalies in financial transactions \cite{hastie2009}.

Reinforcement learning is built around trial and error instead of examples. An agent acts
inside an environment and gets a reward signal telling it how good or bad that action
was, without ever being told what the right action would have been \cite{sutton2018}.
That is the key difference from supervised learning: the feedback only evaluates what
happened, it does not instruct. The agent constantly weighs exploiting what already works
against exploring something that might work better, formalized through the Markov
decision process, where the next state depends only on the current one and the action
taken. Foundational methods include dynamic programming, temporal-difference learning,
and Q-learning, and pairing these ideas with deep networks produced algorithms like Deep
Q-Networks \cite{sutton2018}. It is the paradigm behind game-playing systems such as
AlphaGo, as well as robotic control and autonomous vehicle navigation.

None of this changes with ensemble learning: what changes is not the kind of feedback a
model receives, but what is done with several trained models once they already exist.

\subsection*{Ensemble learning}

A random forest does not introduce a new kind of training signal: it is still a
supervised model. Ensemble learning trains several base models, often called weak
learners, and combines their outputs into a single prediction that in practice tends to
beat any of the individual models it is built from \cite{mienye2022,mohammed2023}. A high
score from one ensemble does not by itself prove that its training procedure generalizes
better than a single model's, and a credible comparison still needs a clearly defined
objective, a proper train/test separation, and metrics that match the question being
asked. The logic behind why ensembles tend to help at all is close to
the wisdom of crowds: if each base model does at least a bit better than random guessing,
and their mistakes are not all correlated, combining their votes cancels out a good part
of the individual errors. The bias-variance decomposition gives this a more formal
footing. A model can go wrong by being too simple for the problem (high bias) or too
sensitive to the particular training sample it saw (high variance); averaging several
independently trained models lowers variance without raising bias, while reweighting
examples so later models correct earlier mistakes instead lowers bias. Formally, if $M$
predictors have equal error variance $\sigma^2$ and equal pairwise error correlation
$\rho$, the variance of their average is
\begin{equation}
  \mathrm{Var}\left(\frac{1}{M}\sum_{m=1}^{M} \hat{f}_m(x)\right)
  = \sigma^2 \left(\rho + \frac{1-\rho}{M}\right).
  \label{eq:ensemble-variance}
\end{equation}
More models help most when their errors are weakly correlated; they cannot remove an
error shared by every model. This is why both individual accuracy and diversity matter in
a random forest. Breiman's analysis puts an exact number on that trade-off: for a forest
whose trees have mean pairwise correlation $\bar\rho$ and strength $s$ (roughly, how much
better than chance each tree predicts on average), the generalization error is bounded by
\begin{equation}
  PE^* \;\leq\; \bar\rho\,\frac{1-s^2}{s^2}.
  \label{eq:breiman-bound}
\end{equation}
The bound shrinks as correlation $\bar\rho$ goes down or strength $s$ goes up, so a good
ensemble needs base models that are both accurate and diverse in their mistakes
\cite{breiman2001}. Equations~\ref{eq:ensemble-variance} and~\ref{eq:breiman-bound} are
intuitions under their stated assumptions, not a guarantee that any ensemble will improve
on a single well-chosen baseline.

Supervised, unsupervised, and reinforcement learning are defined by what kind of feedback
is available during training: labeled targets, unlabeled structure, or a reward. Ensemble
learning is defined by something else, by what is done with a set of models once they
already exist, and it can be applied on top of any of the three \cite{mienye2022}. Most
ensemble methods were developed for supervised learning, where the overlap is strongest,
but the same idea shows up elsewhere: cluster ensembles combine several unsupervised
partitions into one consensus result, and recent work uses multiple value or policy
estimates together in reinforcement learning to stabilize training \cite{mohammed2023}. It
makes more sense to think of ensemble learning as something that cuts across the three
classical paradigms than as a fourth one next to them.

Three strategies cover most of what ensemble learning does, and they also supply its
representative algorithms, summarized in Table~\ref{tab:strategies}.

\begin{table}[h]
\centering
\caption{Core ensemble strategies and their representative algorithms}
\label{tab:strategies}
\begin{tabularx}{\textwidth}{l X X}
\toprule
\textbf{Strategy} & \textbf{How predictions are combined} & \textbf{Main consideration} \\
\midrule
Bagging &
Train each base model on a different bootstrap sample; combine by voting or averaging.
Random Forest also samples features at each split \cite{breiman2001}. &
Trains in parallel; mainly reduces variance, which is why it works well with unstable
models like deep decision trees. Correlated trees limit the gain. \\
\addlinespace
Boosting &
Fit models sequentially so each one corrects the previous ones' mistakes. AdaBoost
reweights misclassified cases \cite{freund1997}; gradient boosting, XGBoost, and LightGBM
fit each new model to the error gradient instead \cite{hastie2009,mienye2022}. &
Trains sequentially; mainly reduces bias. More exposed to noisy or mislabeled data than
bagging. \\
\addlinespace
Stacking &
Train diverse base learners and feed their predictions into a second-level model that
learns how to combine them \cite{mohammed2023}. &
Can beat a fixed voting rule, but adds complexity; base models' predictions must be
properly held out to avoid leakage. \\
\bottomrule
\end{tabularx}
\end{table}

Across all three strategies, the relevant hyperparameters (number of estimators, tree
depth or feature sampling, boosting learning rate, meta-learner complexity) should be
tuned on training/validation data rather than on the final test set.

These methods show up constantly in applied work. Tree-based ensembles reach over 95\%
accuracy on standard heart disease datasets, and XGBoost models trained on routine blood
tests have flagged likely COVID-19 cases with over 99\% accuracy \cite{mienye2022}. Random
forests combined with feature selection reach similar numbers on credit card fraud
detection, and gradient-boosted trees are a common choice for real estate valuation and
fraud detection more broadly \cite{mohammed2023}.

\subsection*{Connection to Deep Learning}

A neural network represents a parametric function $f(x;\theta)$, whose weights and biases
are learned by backpropagation and gradient-based optimization. Its hidden layers extract
features, while a final layer maps those features to the output. Sufficiently expressive
networks can approximate a broad range of functions, but expressiveness alone does not
ensure successful training or generalization. Ensemble learning can combine independently
trained networks, for example by averaging their predicted class probabilities; the added
computation is worthwhile only if it improves the result over well-chosen individual
models \cite{mohammed2023}. Other learning techniques act at different points in this
process: transfer learning reuses a trained feature-extraction backbone for a new task,
and self-supervised learning can train such a backbone with a pretext task that does not
require manual labels. These methods can be combined with an ensemble, but combining
models is not the same as learning a transferable representation; their relationship will
be revisited after the other groups' class presentations, once all the learning
techniques covered this term can be compared side by side.

\subsection*{A preliminary experiment}

As a first practical check, our group implemented and compared a single decision tree
against several ensemble methods (random forest, AdaBoost, gradient boosting, and a
stacking ensemble) on a real classification task, evaluated with accuracy, sensitivity,
specificity, balanced accuracy, and ROC-AUC. Two questions guided the experiment: whether
the ensembles actually outperform a single model past a trivial baseline, and whether the
model ranked best by a single metric is also the model that is easiest to operate safely
once a decision threshold has to be chosen. The general finding is that a single top
metric, taken alone, does not tell the whole story: the model with the highest score on
one criterion is not necessarily the one that reaches a safe operating point with the
least compromise, which is why the relative cost of a missed positive case versus a false
alarm has to guide that choice explicitly rather than being left at a default threshold.
The full experimental setup, results tables, and threshold analysis are reported in the
companion MLLB report.

\subsection*{Conclusion and notes on sources}

Ensembling is most useful when several reasonable models make different errors. Bagging
chiefly addresses instability, boosting learns sequential corrections, and stacking
learns a combination from held-out predictions. None of the three removes the need for a
sensible baseline, a sound train/test separation, and metrics matched to the application.
This appendix reflects the group's first assignment on Ensemble Learning; later versions
of this report will connect it to the learning techniques covered by the other groups and
to the Deep Learning use case developed in the main sections above.

\medskip
\noindent\textit{Note on sources.} Breiman's 1996 paper on bagging and Wolpert's 1992
paper on stacking are foundational to this topic but were not read directly; they are
known to us through Mienye and Sun's 2022 survey \cite{mienye2022}, which is cited here
instead wherever those ideas are discussed.

\end{document}
